\subsection{Orthonormal Coordinates and Inner-Product Identities}

\begin{theorem} \label{thm:orthonormal-basis-coordinates} \hlt{Coordinates in an Orthonormal Basis}\\
Let $\mathcal{V}$ be finite-dimensional and let $\beta=(\mathbf{e}_1,\ldots,\mathbf{e}_n)$ be an orthonormal basis of $\mathcal{V}$.\\
Then for every $\mathbf{v}\in\mathcal{V}$, $\mathbf{v}=\sum\limits_{j=1}^{n}\innerprod{\mathbf{v},\mathbf{e}_j}\mathbf{e}_j$, $\norm{\mathbf{v}}^2=\sum\limits_{j=1}^{n}\abs{\innerprod{\mathbf{v},\mathbf{e}_j}}^2$ and $[\mathbf{v}]_{\beta}=
\begin{bmatrix}
\innerprod{\mathbf{v},\mathbf{e}_1}\\
\vdots\\
\innerprod{\mathbf{v},\mathbf{e}_n}
\end{bmatrix}$.
\begin{proof}
Since $\beta$ is a basis, write $\mathbf{v}=\sum\limits_{j=1}^{n}a_j\mathbf{e}_j$. Taking inner product with $\mathbf{e}_k$ gives $\innerprod{\mathbf{v},\mathbf{e}_k}=\sum\limits_{j=1}^{n}a_j\innerprod{\mathbf{e}_j,\mathbf{e}_k}=a_k$.\\
Thus the displayed expansion and coordinate column follow.\\
Applying Proposition \ref{prop:orthonormal-linear-combination-norm} to the same expansion gives $\norm{\mathbf{v}}^2=\sum\limits_{j=1}^{n}\abs{a_j}^2=\sum\limits_{j=1}^{n}\abs{\innerprod{\mathbf{v},\mathbf{e}_j}}^2$.\\
\hlt{Structural Commentary}: An orthonormal basis makes coordinate extraction an inner-product computation.\\
\hlt{Key Insight}: Pairing with $\mathbf{e}_k$ isolates the $k$th coordinate.
\end{proof}
\end{theorem}

\begin{figure}[H]
  \centering
  \begin{tikzpicture}[
  >=Stealth,
  every node/.style={font=\small},
  axis/.style={gray!60, line width=0.4pt, ->},
  vec/.style={->, blue!75!black, line width=0.95pt},
  comp/.style={->, teal!70!black, line width=0.85pt},
  guide/.style={gray!45, dashed, line width=0.45pt},
  point/.style={circle, fill=black, inner sep=1pt}
]
  \coordinate (O) at (0,0);
  \coordinate (E1) at (3.2,0);
  \coordinate (E2) at (0,2.25);
  \coordinate (V) at (2.45,1.55);
  \coordinate (X) at (2.45,0);
  \coordinate (Y) at (0,1.55);

  \draw[axis] (-0.25,0) -- (3.65,0) node[right] {$\mathbf{e}_1$};
  \draw[axis] (0,-0.25) -- (0,2.65) node[above] {$\mathbf{e}_2$};
  \draw[guide] (X) -- (V) -- (Y);
  \draw[comp] (O) -- (X) node[midway, below] {$\innerprod{\mathbf{v},\mathbf{e}_1}\mathbf{e}_1$};
  \draw[comp] (X) -- (V) node[midway, right] {$\innerprod{\mathbf{v},\mathbf{e}_2}\mathbf{e}_2$};
  \draw[vec] (O) -- (V) node[midway, above left] {$\mathbf{v}$};

  \node[point, label=below left:{$\mathbf{0}$}] at (O) {};
  \node[point] at (V) {};
\end{tikzpicture}

  \caption{An orthonormal basis resolves a vector into perpendicular coordinate components.}
  \label{fig:orthonormal-coordinates}
\end{figure}

\begin{remark} \label{rem:cross-terms} \hlt{Cross Terms}\\
In an expansion such as $\norm{\mathbf{u}+\mathbf{v}}^2$, the terms involving both vectors, $\innerprod{\mathbf{u},\mathbf{v}}$ and $\innerprod{\mathbf{v},\mathbf{u}}$, are called \hlt{cross terms}. The terms $\norm{\mathbf{u}}^2$ and $\norm{\mathbf{v}}^2$ are the pure squared-length terms.
\end{remark}

\begin{theorem} \label{thm:pythagorean-inner-product} \hlt{Pythagorean Theorem}\\
If $\mathbf{u},\mathbf{v}\in\mathcal{V}$ and $\mathbf{u}\perp\mathbf{v}$, then $\norm{\mathbf{u}+\mathbf{v}}^2=\norm{\mathbf{u}}^2+\norm{\mathbf{v}}^2$.
\begin{proof}
Since $\mathbf{u}\perp\mathbf{v}$, Proposition \ref{prop:orthogonality-and-zero} gives $\innerprod{\mathbf{v},\mathbf{u}}=0$. Thus
\begin{align}
\norm{\mathbf{u}+\mathbf{v}}^2 =\innerprod{\mathbf{u}+\mathbf{v},\mathbf{u}+\mathbf{v}}=\innerprod{\mathbf{u},\mathbf{u}}+\innerprod{\mathbf{u},\mathbf{v}}
  +\innerprod{\mathbf{v},\mathbf{u}}+\innerprod{\mathbf{v},\mathbf{v}} =\norm{\mathbf{u}}^2+\norm{\mathbf{v}}^2.
\nonumber
\end{align}
\hlt{Structural Commentary}: Orthogonality cancels the cross terms in the squared-norm expansion.\\
\hlt{Key Insight}: Perpendicular components add in squared length, not in length.
\end{proof}
\end{theorem}

\begin{proposition} \label{prop:pythagorean-converse-real-inner-products} \hlt{Pythagorean Converse over $\mathbb{R}$}\\
In a real inner-product space, $\mathbf{u}\perp\mathbf{v}$ $\Longleftrightarrow$ $\norm{\mathbf{u}+\mathbf{v}}^2=\norm{\mathbf{u}}^2+\norm{\mathbf{v}}^2$.
\begin{proof}
Use $\norm{\mathbf{u}+\mathbf{v}}^2 =\norm{\mathbf{u}}^2+2\innerprod{\mathbf{u},\mathbf{v}}+\norm{\mathbf{v}}^2$.
\begin{enumerate}[label=\roman*.]
\setlength{\itemsep}{0pt}
\item[$(\Rightarrow)$] If $\mathbf{u}\perp\mathbf{v}$, then $\innerprod{\mathbf{u},\mathbf{v}}=0$, so the identity follows.
\item[$(\Leftarrow)$] If the identity holds, then $2\innerprod{\mathbf{u},\mathbf{v}}=0$, hence $\mathbf{u}\perp\mathbf{v}$.
\end{enumerate}
\hlt{Structural Commentary}: In real inner-product spaces, the cross term has no hidden imaginary part.\\
\hlt{Key Insight}: The Pythagorean identity detects orthogonality exactly over $\mathbb{R}$.
\end{proof}
\end{proposition}

\begin{theorem} \label{thm:bessel-inequality-finite-orthonormal-list} \hlt{Bessel's Inequality}\\
Let $(\mathbf{e}_1,\ldots,\mathbf{e}_m)$ be an orthonormal list in an inner-product space $\mathcal{V}$.\\
Then, for every $\mathbf{v}\in\mathcal{V}$, $\sum\limits_{j=1}^{m}\abs{\innerprod{\mathbf{v},\mathbf{e}_j}}^2\leq\norm{\mathbf{v}}^2$.
\begin{proof}
Set $\mathbf{u}=\sum_{j=1}^{m}\innerprod{\mathbf{v},\mathbf{e}_j}\mathbf{e}_j$. For each $k$, orthonormality gives $\innerprod{\mathbf{v}-\mathbf{u},\mathbf{e}_k} =\innerprod{\mathbf{v},\mathbf{e}_k} -\sum\limits_{j=1}^{m}\innerprod{\mathbf{v},\mathbf{e}_j}\innerprod{\mathbf{e}_j,\mathbf{e}_k}=0$.\\
Thus $\mathbf{v}-\mathbf{u}$ is orthogonal to $\mathbf{u}$. By Pythagorean Theorem \ref{thm:pythagorean-inner-product} and Proposition \ref{prop:orthonormal-linear-combination-norm},
\begin{align}
\norm{\mathbf{v}}^2=\norm{\mathbf{u}}^2+\norm{\mathbf{v}-\mathbf{u}}^2
\geq \norm{\mathbf{u}}^2
=\sum_{j=1}^{m}\abs{\innerprod{\mathbf{v},\mathbf{e}_j}}^2.
\nonumber
\end{align}
\hlt{Structural Commentary}: The chosen directions give an orthogonal part of $\mathbf{v}$; the residual is perpendicular.\\
\hlt{Key Insight}: Partial orthonormal coordinates cannot contain more energy than the whole vector.
\end{proof}
\end{theorem}

\begin{corollary} \label{cor:orthonormal-sequence-coefficients-vanish} \hlt{Coefficients Against an Orthonormal Sequence Vanish}\\
Let $(\mathbf{e}_1,\mathbf{e}_2,\ldots)$ be an orthonormal sequence in an inner-product space $\mathcal{V}$, meaning every finite sublist is orthonormal. Then, for every $\mathbf{v}\in\mathcal{V}$, $\lim\limits_{n\to\infty}\innerprod{\mathbf{v},\mathbf{e}_n}=0$.
\begin{proof}
Fix $\epsilon>0$. By Theorem \ref{thm:bessel-inequality-finite-orthonormal-list}, every finite set of indices satisfying $\abs{\innerprod{\mathbf{v},\mathbf{e}_j}}\geq\epsilon$ has cardinality at most $\norm{\mathbf{v}}^2/\epsilon^2$. Hence only finitely many $j$ have $\abs{\innerprod{\mathbf{v},\mathbf{e}_j}}\geq\epsilon$. This holds for every $\epsilon>0$, so $\innerprod{\mathbf{v},\mathbf{e}_n}\to0$.\\
\hlt{Structural Commentary}: Bessel bounds the number of large orthonormal coefficients.\\
\hlt{Key Insight}: Infinitely many large coefficients would force infinite energy.
\end{proof}
\end{corollary}

\begin{theorem} \label{thm:parseval-identity-finite-dimensional} \hlt{Parseval's Identity}\\
Let $\beta=(\mathbf{e}_1,\ldots,\mathbf{e}_n)$ be an orthonormal basis of an inner-product space $\mathcal{V}$.\\
Then, for all $\mathbf{v},\mathbf{w}\in\mathcal{V}$, $\innerprod{\mathbf{v},\mathbf{w}} =\sum\limits_{j=1}^{n}\innerprod{\mathbf{v},\mathbf{e}_j}\overline{\innerprod{\mathbf{w},\mathbf{e}_j}} =\innerprod[\mathbb{F}^n]{[\mathbf{v}]_{\beta},[\mathbf{w}]_{\beta}}$.
\begin{proof}
By Theorem \ref{thm:orthonormal-basis-coordinates}, $\mathbf{v}=\sum_{j=1}^{n}\innerprod{\mathbf{v},\mathbf{e}_j}\mathbf{e}_j$ and $\mathbf{w}=\sum_{k=1}^{n}\innerprod{\mathbf{w},\mathbf{e}_k}\mathbf{e}_k$. By linearity in first slot and conjugate linearity in second, $\innerprod{\mathbf{v},\mathbf{w}} =\sum\limits_{j=1}^{n}\sum\limits_{k=1}^{n} \innerprod{\mathbf{v},\mathbf{e}_j}\overline{\innerprod{\mathbf{w},\mathbf{e}_k}} \innerprod{\mathbf{e}_j,\mathbf{e}_k} =\sum\limits_{j=1}^{n}\innerprod{\mathbf{v},\mathbf{e}_j}\overline{\innerprod{\mathbf{w},\mathbf{e}_j}}$.\\
The last expression is exactly the standard inner product of the coordinate columns in $\mathbb{F}^n$.\\
\hlt{Structural Commentary}: Orthonormal coordinates make inner products coordinatewise.\\
\hlt{Key Insight}: Orthonormal coordinates preserve angles and lengths, not just linear combinations.
\end{proof}
\end{theorem}

\begin{theorem} \label{thm:matrix-entries-from-orthonormal-codomain} \hlt{Matrix Entries from an Orthonormal Codomain Basis}\\
Let $\mathcal{V}$ and $\mathcal{W}$ be finite-dimensional vector spaces over $\mathbb{F}$, with $\mathcal{W}$ an inner-product space. Let $\beta=(\mathbf{v}_1,\ldots,\mathbf{v}_n)$ be a basis of $\mathcal{V}$ and let $\gamma=(\mathbf{w}_1,\ldots,\mathbf{w}_m)$ be an orthonormal basis of $\mathcal{W}$. If $T\in\LL(\mathcal{V},\mathcal{W})$, then
\begin{align}
[T]_{\gamma\leftarrow\beta}
=\begin{bmatrix}
\innerprod{T\mathbf{v}_1,\mathbf{w}_1} & \cdots & \innerprod{T\mathbf{v}_n,\mathbf{w}_1}\\
\vdots & \ddots & \vdots\\
\innerprod{T\mathbf{v}_1,\mathbf{w}_m} & \cdots & \innerprod{T\mathbf{v}_n,\mathbf{w}_m}
\end{bmatrix}.
\nonumber
\end{align}
Conversely, for each $A\in\operatorname{Mat}(m,n,\mathbb{F})$, there is a unique $T\in\LL(\mathcal{V},\mathcal{W})$ such that $[T]_{\gamma\leftarrow\beta}=A$.
\begin{proof}
The $j$th column of $[T]_{\gamma\leftarrow\beta}$ is $[T\mathbf{v}_j]_{\gamma}$. Since $\gamma$ is orthonormal, Theorem \ref{thm:orthonormal-basis-coordinates} gives
\begin{align}
[T\mathbf{v}_j]_{\gamma}
=\begin{bmatrix}
\innerprod{T\mathbf{v}_j,\mathbf{w}_1}\\
\vdots\\
\innerprod{T\mathbf{v}_j,\mathbf{w}_m}
\end{bmatrix}.
\nonumber
\end{align}
This proves the displayed matrix formula. The converse is exactly the surjectivity and injectivity of the coordinate encoding in Theorem \ref{thm:linear-maps-and-matrices-same-data}.\\
\hlt{Structural Commentary}: Columns are basis-image coordinates, read by inner products.\\
\hlt{Key Insight}: Each entry measures $T\mathbf{v}_j$ against one output direction.
\end{proof}
\end{theorem}

\begin{figure}[H]
  \centering
  \begin{tikzpicture}[
  x=0.88cm,
  y=0.88cm,
  >=Stealth,
  every node/.style={font=\small},
  space/.style={draw=gray!45, rounded corners=3pt, minimum width=3.25cm, minimum height=2.3cm},
  vec/.style={->, blue!75!black, line width=0.9pt},
  basis/.style={->, teal!70!black, line width=0.8pt},
  map/.style={->, gray!70, line width=0.7pt},
  measure/.style={->, orange!80!black, line width=0.65pt}
]
  \node[space] (V) at (0,0) {};
  \node[anchor=north west] at (-1.65,1.15) {$\mathcal{V}$};
  \node[space] (W) at (5.4,0) {};
  \node[anchor=north west] at (3.75,1.15) {$\mathcal{W}$};

  \coordinate (Ov) at (-0.95,-0.55);
  \coordinate (vj) at (0.85,0.65);
  \draw[vec] (Ov) -- (vj) node[midway, above left] {$\mathbf{v}_j$};
  \node[below] at (Ov) {$0$};

  \coordinate (Ow) at (4.55,-0.75);
  \coordinate (Tvj) at (6.45,0.75);
  \coordinate (wone) at (6.1,-0.75);
  \coordinate (wtwo) at (4.55,0.75);
  \draw[basis] (Ow) -- (wone) node[right] {$\mathbf{w}_1$};
  \draw[basis] (Ow) -- (wtwo) node[above] {$\mathbf{w}_2$};
  \draw[vec] (Ow) -- (Tvj) node[midway, above left] {$T\mathbf{v}_j$};
  \draw[gray!45, dashed] (Tvj) -- (6.45,-0.75);
  \draw[gray!45, dashed] (Tvj) -- (4.55,0.75);

  \draw[map] (1.95,0) -- (3.35,0) node[midway, above] {$T$};

  \node[draw=gray!45, rounded corners=3pt, align=center, inner sep=5pt] (coords) at (9.25,-0.05)
    {$[T\mathbf{v}_j]_{\gamma}$\\[2pt]
     $=\begin{bmatrix}
        \innerprod{T\mathbf{v}_j,\mathbf{w}_1}\\
        \vdots\\
        \innerprod{T\mathbf{v}_j,\mathbf{w}_m}
       \end{bmatrix}$};
  \draw[measure] (7.05,0.35) -- (coords.west);
\end{tikzpicture}

  \caption{The $j$th column of $[T]_{\gamma\leftarrow\beta}$ is obtained by resolving $T\mathbf{v}_j$ against the orthonormal basis of $\mathcal{W}$.}
  \label{fig:orthonormal-matrix-entries}
\end{figure}

\begin{corollary} \label{cor:inner-product-through-operator-matrix} \hlt{Inner Products Through an Operator Matrix}\\
Let $\beta=(\mathbf{e}_1,\ldots,\mathbf{e}_n)$ be an orthonormal basis of a finite-dimensional inner-product space $\mathcal{V}$, and let $T\in\LL(\mathcal{V})$. Then, for all $\mathbf{v},\mathbf{w}\in\mathcal{V}$, $\innerprod{T\mathbf{v},\mathbf{w}} =\innerprod[\mathbb{F}^n]{[T]_{\beta\leftarrow\beta}[\mathbf{v}]_{\beta},[\mathbf{w}]_{\beta}}$.
\begin{proof}
By Proposition \ref{prop:matrix-representation-action}, $[T\mathbf{v}]_{\beta}=[T]_{\beta\leftarrow\beta}[\mathbf{v}]_{\beta}$. Applying Theorem \ref{thm:parseval-identity-finite-dimensional} to $T\mathbf{v}$ and $\mathbf{w}$ gives the result.\\
\hlt{Structural Commentary}: Orthonormal coordinates turn $\innerprod{T\mathbf{v},\mathbf{w}}$ into matrix computation.\\
\hlt{Key Insight}: Orthonormal coordinates make operator geometry computable by matrices.
\end{proof}
\end{corollary}
